\section{Conclusiones}

El trabajo permitió analizar el tipo de cambio UYU/BRL como una variable cuantitativa continua observada en el tiempo. La muestra contiene \TotalRegistros{} registros entre \FechaMinima{} y \FechaMaxima{}. La media muestral fue \MediaCambio{}, la mediana fue \MedianaCambio{} y el desvío estándar fue \DesvioCambio{}, lo que resume el nivel central y la variabilidad de la cotización en el período estudiado.

La distribución empírica fue explorada mediante histograma, boxplot, cuartiles e IQR. No se detectaron outliers por el criterio de 1.5 IQR, pero sí se observaron cambios de nivel y variabilidad temporal. Por ello, la variable no debe tratarse como una muestra independiente sin estructura temporal.

Entre las distribuciones candidatas evaluadas para el nivel de la cotización, la Log-normal presentó el mejor ajuste relativo según KS y chi-cuadrado. Sin embargo, la evidencia no permite afirmar que esa distribución describa exactamente el fenómeno real. El ajuste debe interpretarse como una aproximación útil, no como una verdad probabilística definitiva.

Los intervalos de confianza para la media y la varianza aplican herramientas centrales del curso, pero su interpretación debe considerar la dependencia temporal. Para mejorar el análisis se recomienda incorporar pruebas de autocorrelación conjunta, bootstrap por bloques, modelos específicos de series temporales, análisis por subperíodos y comparación con otras fuentes de datos.

En síntesis, el análisis describe el comportamiento observado del tipo de cambio UYU/BRL, cuantifica su variabilidad y evalúa modelos probabilísticos posibles. Su utilidad principal es organizar evidencia estadística sobre el período estudiado, no producir una predicción cierta ni establecer relaciones causales.

